% DRAFT Introduction (Sprint 1). Owner: Trung.
% Drop-in replacement for sections/01_introduction.tex in the IEEE scaffold
% (branch origin/trungnguyen). New BibTeX keys are in drafts/references-additions.bib;
% merge them into references.bib together with this section.
%
% Rules followed: every number below is a measured Sprint 1 observation
% (docs/sprint1/report.md). Nothing about FinObs results is claimed yet;
% contributions are stated as goals. Lines marked TODO need team confirmation.

\section{Introduction}
\label{sec:introduction}

Financial decisions increasingly rest on data that is collected, joined, and
transformed by long chains of automated jobs. A trading dashboard, a risk report,
or a model feature for Vietnamese equities may depend on tick streams from several
exchanges, an ingestion service, a message queue, streaming and batch
transformations, and a lakehouse table format~\cite{armbrust2021lakehouse}.
Each hop can corrupt the data without failing: a job that finishes
successfully on stale, duplicated, or mis-typed input still produces a
plausible-looking number. Such \emph{silent data failures} are particularly costly
in finance, where downstream consumers act on the result within seconds and
rarely have the context to question it.

Traditional monitoring watches the infrastructure (CPU, memory, job status), not
the data. Data validation libraries close part of this gap by checking declared
constraints on each batch~\cite{schelter2018automating}, but they assume that the
engineer already knows which constraints matter. In practice, many incidents
violate no single rule. They appear only as a change in freshness, volume, or
distribution, or as a disagreement between sources that should agree. Even once an
anomaly is detected, finding \emph{where} it originated in a multi-stage pipeline
remains a largely manual task that requires knowledge of the lineage between
datasets.

\subsection{Why financial market data is a hard case}

Market data combines several properties that make reliability difficult.
First, it is unbounded and arrives out of order, so correctness depends on
event-time semantics and watermarks rather than arrival order~\cite{akidau2015dataflow}.
Second, the same asset is quoted by multiple independent sources, and those
sources do not always agree: Li et al.~\cite{li2012truth} found substantial
inconsistency even for seemingly objective values such as stock prices.
Third, sources differ in protocol and semantics: push-based WebSocket feeds for
cryptocurrency exchanges versus a rate-limited polling API for the Ho Chi Minh
Stock Exchange (HOSE), with different time formats, precision, and trading hours.

A one-hour recording collected in the first sprint of this project
(231\,859 trades from Binance and Coinbase over 20 symbols) illustrates how quickly
these issues surface even without any injected fault:
\begin{itemize}
  \item \textbf{Clock skew at the collector.} The collector clock was initially
        1.8\,s off and kept drifting after correction. As a result,
        32\,461 Binance trades had a \emph{negative} measured ingestion delay.
  \item \textbf{Mid-session history replay.} Without any reconnect, Coinbase
        re-sent historical trades as ordinary updates, some of them 19.7 hours old.
  \item \textbf{Duplicates under load.} During a burst, Coinbase delivered the same
        trades more than once.
  \item \textbf{Correlated bursts.} Both exchanges peaked in the same second, at
        43 times Binance's mean rate. This is a genuine market event that a naive
        volume detector would flag as an incident.
\end{itemize}
None of these events caused a pipeline failure, and none would be caught by
infrastructure monitoring. Each one, however, would distort a downstream latency
SLO, a traded-volume aggregate, or a cross-exchange comparison.

\subsection{Root causes are mostly about data, not code}

The types of failures above are consistent with empirical studies of data
pipelines. Foidl et al.~\cite{foidl2023pipeline} analysed developer issues and
questions about data pipelines. They report that problems with data types are the
most frequent root cause of data-related issues, and that ingestion, integration,
and cleaning are the stages where developers struggle most.
% TODO: verify the exact wording and percentages against the paper before citing numbers.
This suggests that an observability system for financial data should be organised
around a taxonomy of \emph{data} faults: stale feeds, partial loss, duplicates,
replayed history, late data, clock skew, schema drift, type and precision errors,
value anomalies, cross-source mismatch, and job failures. Each fault should be
linked to the signals it produces and to the pipeline stage where it originates.

\subsection{Research questions}

FinObs investigates whether detection and diagnosis of such faults can be
automated and evaluated reproducibly. We address three research questions:
\begin{description}
  \item[RQ1 (Detection).] How quickly and how precisely can per-source freshness,
        volume, distribution, schema, and pipeline-health signals detect injected
        data faults in a multi-source financial pipeline, measured by mean time to
        detect (MTTD), precision, and recall?
  \item[RQ2 (Diagnosis).] Can an agent that queries metrics, data contracts, and
        dataset lineage locate the root-cause asset of an incident? We measure this
        with top-$k$ accuracy (AC@$k$) and mean time to diagnose (MTTDiag).
  \item[RQ3 (Cross-source evidence).] Does comparing independent sources of the
        same instrument reduce false alarms on genuine market events and reveal
        faults that single-source checks miss?
\end{description}

\subsection{Contributions}

This work aims to make the following contributions.
% TODO (team): keep only the items that are delivered by the final sprint.
\begin{enumerate}
  \item A \textbf{fault taxonomy} for multi-source financial data pipelines with
        eleven fault types. It is grounded in the root-cause categories of Foidl
        et al.~\cite{foidl2023pipeline} and in faults observed in real exchange data.
  \item A \textbf{labelled scenario benchmark}: a YAML format that separates fault
        injection from ground truth. The ground truth includes the root-cause
        lineage asset, the expected signals, and signals that must \emph{not}
        alert. The benchmark also contains negative-control scenarios built from
        real market events.
  \item An \textbf{observability layer} over a lakehouse pipeline that ingests
        Binance, Coinbase, and HOSE VN30 data. It includes explicit data contracts
        and per-source watermarks.
  \item A \textbf{diagnosis agent} that uses tools over metrics and lineage to
        propose ranked root causes, evaluated against the benchmark ground truth.
\end{enumerate}

\subsection{Paper organisation}

Section~\ref{sec:related} reviews related work on data validation, stream
processing, and truth discovery. Section~\ref{sec:problem} defines the problem and
the requirements. Section~\ref{sec:architecture} describes the FinObs architecture,
and Section~\ref{sec:methodology} the detection and diagnosis methods.
Section~\ref{sec:evaluation} presents the evaluation protocol and
Section~\ref{sec:results} the results. Section~\ref{sec:conclusion} concludes.
